\documentclass[12pt]{article}
\usepackage{../../ficha}

\begin{document}

\ficha{Logaritmos, Equações e Modelos}

\instrucoes{Resolve os exercícios. Lembra-te: $\log_a b = c \Leftrightarrow a^c = b$, com $a > 0$, $a \neq 1$ e $b > 0$. $\ln$ representa o logaritmo de base $e$.}

\secao{Definição de Logaritmo}{Calcula.}
\begin{exercicios}
  \item $\log_2 8$
  \item $\log_3 81$
  \item $\log_5 \dfrac{1}{25}$
  \item $\log_{10} 0{,}001$
  \item $\log_4 2$
  \item $\log_{\frac{1}{2}} 16$
  \item $\ln e^5$
  \item $\ln 1$
  \item $e^{\ln 7}$
  \item $\log_9 27$
\end{exercicios}

\secao{Propriedades dos Logaritmos}{Simplifica, usando as propriedades.}
\begin{exercicios}
  \item $\log_2 12 - \log_2 3$
  \item $\log_6 4 + \log_6 9$
  \item $3\log_5 2 + \log_5 \dfrac{25}{8}$
  \item $\ln(e^2 \cdot \sqrt{e})$
  \item $\log_3 \sqrt[4]{27}$
  \item $\ln 8 - 3\ln 2$
  \item Escreve $\ln\dfrac{x^2 y}{z^3}$ em função de $\ln x$, $\ln y$ e $\ln z$.
  \item Escreve $2\ln a - \dfrac{1}{2}\ln b + \ln c$ como um único logaritmo.
  \item Sabendo que $\log_2 3 \approx 1{,}585$, calcula $\log_2 18$ e $\log_2 \dfrac{9}{4}$.
  \item Usa a mudança de base para escrever $\log_5 20$ em função do $\ln$.
\end{exercicios}

\secao{A Função Logarítmica}{Resolve.}
\begin{exercicios}
  \item Indica o domínio, o contradomínio, a assíntota e a monotonia de $f(x) = \ln x$.
  \item Determina o domínio de $f(x) = \ln(x - 3)$.
  \item Determina o domínio de $g(x) = \log_2(4 - x^2)$.
  \item Determina o domínio de $h(x) = \ln\left(\dfrac{x + 1}{x - 2}\right)$.
  \item Determina a inversa de $f(x) = 2 + \ln(x - 1)$.
  \item Determina a inversa de $g(x) = 3e^{2x} - 1$.
  \item Os gráficos de $y = e^x$ e $y = \ln x$ são simétricos em relação a que reta?
\end{exercicios}

\secao{Equações Logarítmicas e Exponenciais}{Resolve em $\R$.}
\begin{exercicios}
  \item $\log_3 x = 4$
  \item $\ln(2x - 1) = 0$
  \item $\log_2(x + 3) = 5$
  \item $e^x = 5$
  \item $3^{x} = 20$
  \item $\ln x + \ln(x - 2) = \ln 8$
  \item $\log_2 x + \log_2(x + 6) = 4$
  \item $(\ln x)^2 - 3\ln x + 2 = 0$
  \item $e^{2x} - 5e^x + 6 = 0$
  \item $\log_x 49 = 2$
\end{exercicios}

\secao{Inequações Logarítmicas}{Resolve em $\R$.}
\begin{exercicios}
  \item $\ln x > 2$
  \item $\log_2(x - 1) \le 3$
  \item $\log_{\frac{1}{3}} x > -2$
  \item $\ln(x + 2) < \ln(3x - 4)$
  \item $e^{x} > 10$
\end{exercicios}

\secao{Modelos}{Resolve. Podes usar calculadora; arredonda às centésimas.}
\begin{exercicios}
  \item Um capital de 5\,000 euros é aplicado à taxa anual de 4\,\% em juro composto. Ao fim de quanto tempo atinge 8\,000 euros?
  \item A massa de uma substância radioativa é $m(t) = 100e^{-0{,}05t}$ (gramas, $t$ em anos). Determina o tempo de semivida.
  \item O pH de uma solução é $\text{pH} = -\log_{10}[\text{H}^+]$. Determina o pH de uma solução com $[\text{H}^+] = 3{,}2 \times 10^{-4}$.
  \item A magnitude de um sismo na escala de Richter é $M = \log_{10}\dfrac{I}{I_0}$. Quantas vezes mais intenso é um sismo de magnitude 7 do que um de magnitude 5?
  \item Uma população de bactérias duplica a cada 3 horas. Partindo de 1\,000 bactérias, ao fim de quanto tempo haverá 50\,000?
\end{exercicios}

\secao{Desafios}{Resolve os problemas.}
\begin{exercicios}
  \item Resolve: $\log_2(x + 1) - \log_2(x - 1) = 1$.
  \item Mostra que $\log_a b \cdot \log_b a = 1$.
  \item Resolve: $x^{\ln x} = e^4$.
  \item Resolve: $2^{x} = 3^{x - 1}$, apresentando a solução na forma $\dfrac{\ln a}{\ln b}$.
\end{exercicios}

\end{document}
